\section{Review of Related Literature}

\subsection{Economic Impact and Post-Harvest Processing of Mangoes}
Mangoes represent a vital component of the global agricultural trade, and their
economic impact in the Philippines is profound \cite{fernandez2017mango}. Proper
post-harvest processing and quality assessment are essential for maintaining this
economic value. Typically, mangoes are graded and sorted based on their size,
weight, and morphological attributes \cite{Villaruz2025AVA}. In traditional
agricultural settings, this is done manually, which is not only time-consuming
but also subjective and prone to inconsistencies \cite{amdpwe2023}. While physical
weighing scales and volume measurement techniques (such as water displacement)
are accurate, they are slow and occasionally destructive. The development of
automatic visual inspection systems is thus a highly sought-after intervention to
ensure consistency in size across fruit batches and to meet strict market demands.

\subsection{Properties of the Indian Mango (Katchamitha)}
The Indian mango, locally known as Katchamitha \cite{Katchamitha2025}, is the fourth
most produced mango variant in the Philippines. It accounts for a production volume
of 26,268 metric tons, occupying approximately 3.4\% of the total 763,298 metric
tons of national mango volume projected for 2025 \cite{mangoPhilippineOpenStats}.
Despite its popularity as a locally consumed fruit, it presents unique challenges
for automated grading. Katchamitha mangoes exhibit a deceptively green skin even
upon reaching full physiological maturity, sometimes only displaying a faint
yellow blush \cite{Katchamitha20252}. Consequently, ripeness cannot be reliably
determined through conventional RGB color analysis alone. However, as mangoes
ripen, their internal chemical composition changes, which directly influences
their density \cite{fruitDensity}. Thus, estimating the density and volume
accurately provides a more reliable indicator of the fruit's sensory quality and
ripeness phase.

\subsection{Vision-Based Morphological Estimation}
To automate fruit grading, various vision-based techniques have been proposed. 
Sa'ad \textit{et al.} \cite{shape2015mango} utilized visible imaging combined
with a Fourier-descriptor method to classify mango shapes with high accuracy. 
They further estimated the weight of mangoes by approximating the fruit as a
cylinder, achieving a 94.0\% correlation ($R^2$) between estimated and actual
weights using the equation $w = 2.256V - 157.7$. Similarly, Mon and Aung
\cite{volume2020mango} proposed an image processing algorithm to reconstruct the
3D shape of a mango from a 2D top-view image by analyzing light intensity
distribution to estimate thickness, achieving a volume estimation accuracy of
96.8\%. Aung \textit{et al.} \cite{japan2020mango} also demonstrated weight
estimation from single visible fruit surfaces.

While these 2D-based methods represent significant progress \cite{sweetpotato2022},
they inherently rely on mathematical assumptions about the fruit's symmetry (e.g.,cylindrical or ellipsoidal approximation), which can fail to capture the complex,
asymmetric morphology of actual mangoes.

\subsection{Monocular Depth Estimation and Instance Segmentation}
Recent advancements in deep learning have introduced Monocular Depth Estimation
(MDE) as a powerful tool to bridge the gap between 2D imaging and 3D perception.
MDE predicts a depth map from a single RGB image, offering a low-cost and highly
adaptable alternative to expensive hardware like LiDAR or stereo cameras
\cite{depthanythingv2}. 

For instance, Shang \textit{et al.} \cite{depthcl} highlighted the difficulty
of green-to-green instance segmentation---a common issue when green fruits blend
with the canopy. To address this, they proposed DepthCL-Seg, a framework that
fused RGB and monocular depth data. Their approach demonstrated that
incorporating depth information significantly enhanced the feature representation
of low-contrast target regions, surpassing classical methods like Mask R-CNN.
State-of-the-art MDE models, such as Depth Anything V2 \cite{depth_anything_v2},
provide robust, high-resolution depth maps that can be effectively utilized to
refine bounding boxes and segmentation masks.

However, a fundamental limitation of monocular depth estimation models is their
scale ambiguity: they predict relative depth or scale-variant metric outputs that
fluctuate under different distances, camera zooms, or lighting setups
\cite{depthanythingv2}. To convert these relative depth maps into absolute
physical measurements (e.g., in centimeters or grams), researchers often employ
fiducial markers. Fiducial markers, such as ArUco markers, provide a known
geometric reference in the scene. By detecting the ArUco marker's pixel
dimensions and comparing them with its known physical dimensions, a scaling
factor can be computed to calibrate the depth map and reconstruct accurate 3D
point clouds. 

This study builds upon these findings by proposing two architectures that pair
the cutting-edge YOLOv11 instance segmentation with Depth Anything V2: one
uncalibrated, which leverages scale-invariant ratios to infer morphological
features, and one calibrated, which integrates an ArUco marker to resolve scale
ambiguity. By combining RGB color space with calibrated 3D depth data, the
proposed framework captures the true curvature and thickness of the Katchamitha
mango, moving beyond simple 2D assumptions to accurately estimate weight,
density, and internal stone-to-flesh ratios.
